\emph{(i) The form.} By Fact \claimref{synthesis.imported-sakharov-induced-gravity-coefficients}(a), (b), a field of weight $w$ (in units of one minimally coupled real scalar: $1 - 6\xi_{R}$ for a real scalar of curvature coupling $\xi_{R}$, $1$ for a Weyl fermion, $-4$ for a gauge vector with its ghosts in the Gupta-Bleuler gauge, as in \claimref{synthesis.attachment-modules-give-sakharov-weights-cancelling-at-three-separate-generations}) and mass $m$ has heat trace $(4\pi t)^{-2}e^{-tm^{2}}\bigl(\int r + t\int\mathrm{tr}(E + R/6) + \dots\bigr)$, and with a proper-time cut-off $t \ge \Lambda^{-2}$ its contribution to the coefficient of the induced Einstein-Hilbert term is proportional to $w\int_{\Lambda^{-2}}^{\infty}t^{-2}e^{-tm^{2}}\,dt$; for $m = 0$ this is $w\Lambda^{2}$, which with the normalisation of Fact (b), $-(192\pi^{2}\epsilon)^{-1}B\,R$ at $\epsilon = \Lambda^{-2}$ against $R/(16\pi G)$, gives $1/G_{\mathrm{ind}} = B\Lambda^{2}/(12\pi)$, $B = \sum_{j}w_{j}$. With masses, the substitution $t = u\Lambda^{-2}$ gives $\int_{\Lambda^{-2}}^{\infty}t^{-2}e^{-tm^{2}}dt = \Lambda^{2}E_{2}(x)$, $x = m^{2}/\Lambda^{2}$, $E_{n}(x) = \int_{1}^{\infty}u^{-n}e^{-xu}\,du$. The exponential integrals satisfy $E_{2}' = -E_{1}$, $E_{1}'(x) = -e^{-x}/x$, $E_{2}(0) = 1$, and $E_{1}(x) = -\gamma - \ln x + \sum_{k \ge 1}(-1)^{k+1}x^{k}/(k\,k!)$; the check verifies these relations exactly on the truncated series, with $\ln x$ and $\gamma$ kept symbolic, and reads off $E_{2}(x) = 1 + x(\ln x + \gamma - 1) + O(x^{2}\ln x)$. Hence
\[
\frac{1}{G_{\mathrm{ind}}} = \frac{1}{12\pi}\Bigl[B\Lambda^{2} - \sum_{j}w_{j}m_{j}^{2}\ln\frac{\Lambda^{2}}{m_{j}^{2}} - (1 - \gamma)\sum_{j}w_{j}m_{j}^{2} + O(m^{4}/\Lambda^{2})\Bigr].
\]
The logarithm comes from the $t^{2}$ term of the heat trace, whose part linear in $m^{2}$ and $R$ is the cross term $-m^{2}\,\mathrm{tr}(E + R/6)$ of $e^{-tm^{2}}$ with $t\,\mathrm{tr}(E + R/6)$ (checked on the product of the two series); a $t^{2}$ coefficient produces a logarithm under every regulator with one scale (for dimensional continuation, the pole of the same coefficient), so the coefficient $-\sum w_{j}m_{j}^{2}/(12\pi)$ of $\ln\Lambda^{2}$ does not depend on the regulator, while $B\Lambda^{2}$ and the constant do; in particular a covariant cut-off on the gauge-fixed operators of a vector at gauge parameter $\alpha \neq 1$ changes the vector's weight in $B$.

\emph{Gauge.} All weights below are Feynman-gauge values ($\alpha = 1$, the Gupta-Bleuler form of clause (K1) of the Krein premise): under a covariant cut-off on the gauge-fixed operators a vector's weight is $\alpha - 5$ (\claimref{synthesis.krein-quartet-drops-from-the-induced-einstein-hilbert-coefficient-only-under-a-regulator-that-treats-it-alike}, which owns the weak-scale coefficient $B = 8(1 - 6\xi_{R})$ of the twelve vectors of the forced gauge algebra (Theorem \claimref{standard-model.connection-forced}), forty-eight Weyl fields and the two doublets, used here as an input). For a massive vector the logarithm's coefficient is regulator independent but, at $\xi_{R} \neq 0$, gauge dependent: in the $R_{\alpha}$ gauge the transverse part ($-3$) has mass $m_{V}$, while the longitudinal part (weight $\alpha$, its eigenvalues $(\lambda + \alpha m_{V}^{2})/\alpha$), the Goldstone mode (weight $1 - 6\xi_{R}$) and the ghosts ($-2$) have mass squared $\alpha m_{V}^{2}$, so the coefficient of $m_{V}^{2}\ln\Lambda^{2}$ has weight $-3 + \alpha + (1 - 6\xi_{R})\alpha - 2\alpha = -(3 + 6\xi_{R}\alpha)$ (computed): the Proca value $-3$ for every $\alpha$ at $\xi_{R} = 0$, and $-(3 + \alpha)$ at $\xi_{R} = \tfrac16$, $-4$ at $\alpha = 1$. The gauge dependence at $\xi_{R} \neq 0$ is that of an off-shell quantity: the term $\xi_{R}Rv\,h$ makes the vacuum not stationary at order $R$ (the Nielsen identity, cited from memory).

\emph{(ii) The content at the cut-off.} A field of mass above $\Lambda$ is suppressed by $e^{-tm^{2}}$ for $t \ge \Lambda^{-2}$; a field below it contributes its full weight to $B$ whatever its mass. Under Hypothesis (S), the chain of Theorem \claimref{standard-model.chain} has between $M_{I}$ (about $10^{10}$ to $10^{11}$ GeV) and $M_{U}$ (about $10^{16}$ GeV) the Pati-Salam algebra, of dimension $15 + 3 + 3 = 21$, the three generations of sixteen, the bidoublet $(1,2,2)$ (four real scalars if real, eight if complex), the complex $(10,1,3)$ ($60$ real) and, where named, the complex $(15,2,2)$ ($120$ real). So for a cut-off in that window $B = 48 - 84 + \sum_{j}(1 - 6\xi_{R,j}) = -36 + \sum_{j}(1 - 6\xi_{R,j})$: $-36$ if every scalar is conformally coupled, and $28$, $32$ or $152$ for the three scalar contents if every scalar is minimally coupled (computed). Above $M_{U}$ the store does not fix the content. So the sign and size of $B$ at a cut-off above $M_{I}$ are set by the curvature couplings of the chain's scalars, which no stored result fixes. The chain runs with one Higgs doublet from $M_{Z}$ to $M_{I}$, while the rung field is two elementary doublets (\claimref{synthesis.linear-rung-field-has-a-neutral-normal-pair-in-the-scalon-sum}); this tension between stored results is named here and not resolved, and the scales $M_{I}$, $M_{U}$ quoted are the chain's.

\emph{(iii)(a) The remainder for the weak-scale content alone.} Suppose the content below $\Lambda$ is the weak-scale content and $B = 0$ ($\xi_{R} = \tfrac16$). The scalars have weight $0$; a Dirac fermion of mass $m$ is two Weyl fields of mass $m$; a massive vector has weight $-4$ (the Gauge paragraph, $\alpha = 1$). With $L = \ln(\Lambda^{2}/m_{t}^{2})$ and $\ell_{V} = \ln(m_{t}^{2}/m_{V}^{2}) \le m_{t}^{2}/m_{V}^{2} - 1$,
\[
S = \sum_{j}w_{j}m_{j}^{2}\ln\frac{\Lambda^{2}}{m_{j}^{2}} \;\ge\; \bigl(6m_{t}^{2} - 4(2m_{W}^{2} + m_{Z}^{2})\bigr)L - 4\bigl(2m_{W}^{2}\ell_{W} + m_{Z}^{2}\ell_{Z}\bigr),
\]
since every other fermion adds a positive term when $\Lambda$ exceeds its mass. With the stored values of $m_{t}$, $m_{W}$, $m_{Z}$ the coefficient of $L$ is positive and the constant is less than three times it (computed exactly), and for $\Lambda \ge 10\,m_{t}$, $L \ge \ln 100 > 4$ (as $e < 3$ and $3^{4} < 100$); so $S > 0$ and the loop term $-S/(12\pi)$ is negative. At other $\alpha$ the coefficient of $L$, $6m_{t}^{2} - (3 + \alpha)(2m_{W}^{2} + m_{Z}^{2})$, stays positive up to $\alpha = 5$ (computed). At $\xi_{R} = \tfrac16$ the doublets' classical coupling $-\xi_{R}R\,|\phi|^{2}$ with $|\phi|^{2} = v^{2}/2$ shifts $1/(16\pi G)$ by $-\xi_{R}v^{2}/2$, that is $\Delta(1/G) = -8\pi\xi_{R}v^{2} = -\tfrac{4\pi}{3}v^{2}$ at tree level, of the same sign (computed from the stored $v$; numerically $-2.5 \times 10^{5}$ GeV$^{2}$, about $1.3$ times the loop term at $\Lambda = M_{P}$). For the size at $\Lambda = M_{P}$: $m^{2}\ln(\Lambda^{2}/m^{2})$ increases with $m$ below $\Lambda/\sqrt{e}$, so $|S| \le 48\,m_{t}^{2}\ln(M_{P}^{2}/m_{t}^{2})$ (twenty-four Dirac fermions, the neutrinos included, none heavier than the top; the vectors subtract), and with $\ln y$ below the bit length of $y$ and $\pi > 3$ the check bounds the loop term by $10^{-31}M_{P}^{2}$ and the tree term by $10^{-32}M_{P}^{2}$ (computed); numerically, from the top, $W$ and $Z$ alone, $S/(12\pi M_{P}^{2})$ is about $2 \times 10^{-33}$ (computed, in the proof text only).

\emph{(iii)(b) Under (S).} For a cut-off above $M_{I}$ the remainder contains the fields at $M_{I}$: the $21 - 12 = 9$ vectors that the vev along the Majorana element makes massive, of weight $-4$ each, so $-4 \cdot 9\,M_{I}^{2}\ln(\Lambda^{2}/M_{I}^{2})$ in $S$, of the attractive sign; the chain's scalars near $M_{I}$, whose weights are those that $B = 0$ would need ($36$ units of nonconformal weight) and are not fixed; and, under Premise (M) (\claimref{neutrino.m-framework-realises-majorana-term-scale}), the three right-handed neutrinos, of weight $+1$ and Majorana mass $M$, of the repulsive sign. The fermion argument of (a) does not cover the heavy vectors, so the sign of the remainder is not fixed by the store. Its size is: with total weight at most $4 \cdot 9 + 188 + 3 = 227$ in absolute value (scalar weights of size at most $1$, that is $0 \le \xi_{R,j} \le \tfrac13$), masses up to $6 \times 10^{11}$ GeV (the largest stored $M_{I}$), $\Lambda = M_{P}$ and the same bounds, the check bounds these fields' contribution by $10^{-12}M_{P}^{2}$ (computed). Fields at or above $M_{U}$ are not stored and are not bounded here.

\emph{(iv) Matching.} If $B > 0$, $1/G_{\mathrm{ind}} = M_{P}^{2}$ at leading order exactly when $\Lambda^{2} = 12\pi M_{P}^{2}/B$; for the weak-scale content with minimally coupled doublets ($B = 8$) this is $\Lambda^{2} = \tfrac{3\pi}{2}M_{P}^{2}$, $\Lambda \approx 2.17\,M_{P}$ (computed). The cut-off is then chosen to give $G$, which is a fit, not a prediction.

\emph{Reading (argued).} A prediction of $G$ on this route needs a cut-off scale in Planck units from the series, and the store has none: $N = M_{P}/H$ is initial data (\claimref{programme.n-parameter-of-the-laws}); the weak scale's constant is a second free number (\claimref{synthesis.rung-unit-second-free-number}); the descent's duration $\tau$, which would set the deposit's ultraviolet end $k_{\mathrm{UV}} \sim 1/\tau$, is not fixed, and in any case is a datum of the state, which by the Hadamard property enters the stress tensor only through smooth state-dependent terms, not through the local $G_{ab}$ coefficient that the regulator sets (Fact \claimref{synthesis.imported-microlocal-hadamard-condition}; \claimref{synthesis.hadamard-condition-leaves-the-deposit-spectrum-free-on-bounded-bands}); the floors are a frame-dependent cut-off with no scale in Planck units (\claimref{synthesis.floor-regulator-keeps-the-charge-and-the-central-term-below-the-cut-off}); the Pauli-Villars regulators of the Krein premise replace $\Lambda^{2}$ by $\sum_{i}c_{i}M_{i}^{2}\ln M_{i}^{2}$ with free regulator masses; and the locally covariant scheme has no Einstein-tensor freedom for massless fields and a free coefficient times $m^{2}$ for massive ones (Fact (d)). So every route in the store reaches $G$ only by fixing one dimensionful regulator datum to the observed value, and $G$ stays an input. The mechanism is standard one-loop physics (transport-absent); what is new here is its evaluation on the series' content at and above the weak scale.
